\Needspace{8\baselineskip}
\Article{285}
We can say only a little here about the questions making up our third
problem (to which the first two were reduced above): given two ternary
forms of the same determinant, decide whether they are equivalent, and,
if they are, find all transformations of one into the other. A complete
solution, like the one we gave for the analogous problems for binary
forms, faces still greater difficulties here.
\SourceLeaf{337}
So we will restrict our investigation to certain special cases, for
which this digression was chiefly intended.

I\quad For determinant $+1$, we showed above that all ternary forms
fall into two classes, one containing all indefinite forms, the other
all definite (negative) forms. We immediately conclude that any two
ternary forms of determinant $1$ are equivalent if both are definite
or both indefinite; but if one is definite and the other indefinite,
they are not equivalent (the latter part clearly holds generally for
forms of any determinant). --- Likewise, any two forms of determinant
$-1$ will certainly be equivalent if both are definite or both indefinite.
--- Two definite forms of determinant $2$ will always be equivalent.
Two indefinite ones will not be equivalent if the first three coefficients
of one are all even, but those of the other are not all even; in the
remaining cases (if both have their first three coefficients all even,
or neither does) they will be equivalent. --- We could give many more
special propositions of this kind if we had worked out more examples
above (Article 277).

II.\ In all these cases, if $f$, $f'$ denote equivalent ternary
forms, we can also find one transformation of either into the other.
For, in all cases, we gave above a fairly small set of forms in each
class of ternary forms, to one of which any form of the same class can
be reduced by uniform methods; there we also showed how to reduce all
of these to a single one. Let $F$ be this form in the class containing
$f$, $f'$. The rules given above allow us to find transformations
of the forms $f$, $f'$ into $F$, as well as of $F$ into
$f$, $f'$. From these, Article 270 gives transformations of the
form $f$ into $f'$ and of $f'$ into $f$.

III.\ It would therefore remain only to show how all possible
transformations can be derived from one transformation of the ternary
form $f$ into another $f'$. This problem depends on a simpler
one: finding all transformations of the ternary form $f$ into itself.
For if $f$ is transformed into itself by several substitutions
$(\tau)$, $(\tau')$, $(\tau'')$, and so on, and by the substitution $(t)$
into $f'$, then combining the transformation $(t)$ with $(\tau)$, $(\tau')$,
$(\tau'')$, and so on according to Article 270 clearly gives transformations
that all take $f$ into $f'$. Further, a calculation readily proves
that every transformation of the form $f$ into $f'$ can be derived
in this way by combining the given transformation $(t)$ of the form
$f$ into $f'$ with some (and indeed
\SourceLeaf{338}
\textit{just one}) transformation of the form $f$ into itself.
So combining the given transformation of the form $f$ into $f'$
with \textit{all} transformations of $f$ into itself gives \textit{all}
transformations of $f$ into $f'$, and each exactly once.

\Needspace{11\baselineskip}
Here we restrict the search for all transformations of the form $f$
into itself to the case where $f$ is a definite form whose coefficients
numbered $4,5,6$ are all $=0$\HistoricalNote{*)}{The remaining cases where $f$ is a definite form can be reduced to this one. But if $f$ is an indefinite form, an entirely different method is needed, and there will be infinitely many transformations.}.
So let $f=\left(\begin{smallmatrix}a,\,a',\,a''\\0,\,0,\,0\end{smallmatrix}\right)$, and let all substitutions that take $f$ into itself
be expressed generally by
\[
\begin{array}{ccc}
\alpha,&\beta,&\gamma\\
\alpha',&\beta',&\gamma'\\
\alpha'',&\beta'',&\gamma''
\end{array}
\]
so that the following equations must hold:
\[
\begin{aligned}
a\alpha\alpha+a'\alpha'\alpha'+a''\alpha''\alpha''&=a&&\ldots\quad(\Omega)\\
a\beta\beta+a'\beta'\beta'+a''\beta''\beta''&=a'\\
a\gamma\gamma+a'\gamma'\gamma'+a''\gamma''\gamma''&=a''\\
a\alpha\beta+a'\alpha'\beta'+a''\alpha''\beta''&=0\\
a\alpha\gamma+a'\alpha'\gamma'+a''\alpha''\gamma''&=0\\
a\beta\gamma+a'\beta'\gamma'+a''\beta''\gamma''&=0
\end{aligned}
\]
We now have three cases to distinguish:

I.\ When $a$, $a'$, $a''$ (which will have the same sign) are
all unequal, suppose that $a<a'$, $a'<a''$ (if their order of magnitude
is different, the same conclusions follow in exactly the same way).
Then the first equation in $(\Omega)$ clearly requires $\alpha'=\alpha''=0$, and so
$\alpha=\pm1$. Equations $4,5$ then give $\beta=0$, $\gamma=0$. Likewise, equation
$2$ gives $\beta''=0$, and so $\beta'=\pm1$. From equation $6$ we get
$\gamma'=0$, and from $3$, $\gamma''=\pm1$, so (since the signs can be chosen
independently) there are in all $8$ different transformations.

II.\ When two of the numbers $a$, $a'$, $a''$ are equal,
\textit{for example}, $a'=a''$, and the third is different, suppose

\textit{first} that $a<a'$. Then, as in the preceding case, $\alpha'=0$,
$\alpha''=0$, $\alpha=\pm1$, $\beta=0$, $\gamma=0$. But equations $2,3,6$ readily give
\SourceLeaf{339}
either $\beta'=\pm1$, $\gamma'=0$, $\beta''=0$, $\gamma''=\pm1$, or $\beta'=0$, $\gamma'=\pm1$, $\beta''=\pm1$, $\gamma''=0$.

But if, \textit{second}, $a>a'$, the same conclusions follow in this
way: equations $2,3$ necessarily give $\beta=0$, $\gamma=0$, and either
$\beta'=\pm1$, $\gamma'=0$, $\beta''=0$, $\gamma''=\pm1$, or $\beta'=0$, $\gamma'=\pm1$, $\beta''=\pm1$,
$\gamma''=0$. Under either supposition, equations $4,5$ give $\alpha'=0$,
$\alpha''=0$, and equation $1$ gives $\alpha=\pm1$. So in either case there
are $16$ different transformations. --- The remaining two cases,
where either $a=a''$ or $a=a'$, are handled in exactly the same way,
if we interchange the symbols $\alpha$, $\alpha'$, $\alpha''$ with $\beta$,
$\beta'$, $\beta''$ in the first case, and with $\gamma$, $\gamma'$, $\gamma''$
in the second, respectively.

III.\ When $a$, $a'$, $a''$ are all equal, equations $1,2,3$
require that two of the three numbers $\alpha$, $\alpha'$, $\alpha''$,
and likewise of $\beta$, $\beta'$, $\beta''$, and also of $\gamma$,
$\gamma'$, $\gamma''$, be $=0$, and the third $=\pm1$. But equations
$4,5,6$ readily show that only one of the three numbers $\alpha$,
$\beta$, $\gamma$ can be $=\pm1$, and likewise of $\alpha'$, $\beta'$,
$\gamma'$, and of $\alpha''$, $\beta''$, $\gamma''$. \mbox{So there are only six combinations:}
\[
\begin{array}{c|c|c|c|c|c|c|}
\alpha&\alpha&\alpha'&\alpha'&\alpha''&\alpha''&=\pm1\\
\beta'&\beta''&\beta&\beta''&\beta&\beta'&=\pm1\\
\gamma''&\gamma'&\gamma''&\gamma&\gamma'&\gamma&=\pm1
\end{array}
\quad\text{The other six coefficients }=0
\]
The choices of signs therefore give $48$ transformations in all.
--- The same scheme includes the preceding cases as well: but of the
first six columns we must take only the first when $a$, $a'$,
$a''$ are all unequal; the first and second when $a'=a''$; the first
and third when $a=a'$; the first and sixth when $a=a''$.

It follows that, if the form $f=axx+a'x'x'+a''x''x''$ is transformed into an equivalent
form $f'$ by the substitution
\[
x=\delta y+\varepsilon y'+\zeta y'',\quad x'=\delta'y+\varepsilon'y'+\zeta'y'',\quad x''=\delta''y+\varepsilon''y'+\zeta''y''
\]
all transformations of the form $f$ into $f'$ are contained
in the following scheme:
\[
\begin{array}{c|c|c|c|c|c|l}
x&x&x'&x'&x''&x''&=\pm(\delta y+\varepsilon y'+\zeta y'')\\
x'&x''&x&x''&x&x'&=\pm(\delta'y+\varepsilon'y'+\zeta'y'')\\
x''&x'&x''&x&x'&x&=\pm(\delta''y+\varepsilon''y'+\zeta''y'')
\end{array}
\]
with the distinction that all of the first six columns must be used when $a=a'=a''$;
\SourceLeaf{340}
columns $1$ and $2$ when $a'$, $a''$ are equal and $a$
different; $1$ and $3$ when $a=a'$; $1$ and $6$ when
$a=a''$; finally, only the first column when $a$, $a'$, $a''$
are all unequal. The number of transformations will be $48$ in the
first case, $16$ in the second, third, and fourth, and $8$ in the fifth.

\Needspace{16\baselineskip}
\begin{center}
$\begin{array}{ccc}*&\quad&*\\&*&\end{array}$
\end{center}

\begin{center}
\textit{SOME APPLICATIONS TO THE THEORY OF BINARY FORMS.}
\end{center}

\begin{center}
\textit{Finding a form whose doubling gives a given binary form of the principal genus.}
\end{center}

From this brief account of the first elements of the theory of ternary
forms, we turn to some special applications. First place among them
belongs to the following

\Needspace{8\baselineskip}
\Article{286}
\textsc{Problem.} \textit{Given a binary form $F=(A,B,C)$ of determinant
$D$ belonging to the principal genus, find a binary form $f$
whose doubling gives it.}

\textit{Solution.} I.\ Find a proper representation of the form opposite
to $F$, namely $F'=ATT-2BTU+CUU$, by the ternary form $xx-2yz$, say
\[
x=\alpha T+\beta U,\quad y=\alpha'T+\beta'U,\quad z=\alpha''T+\beta''U
\]
The preceding theory of ternary forms readily shows that this is possible.
For, by hypothesis, $F$ is in the principal genus, so there is a
value of the expression $\surd(A,\allowbreak B,\allowbreak C)(\mathrm{mod.}\,\allowbreak D)$. From it we can find a ternary form
$\varphi$ of determinant $1$ containing $(A,-B,C)$ as a part. All its
coefficients are readily seen to be integers. It is just as easy to see
that $\varphi$ will be an indefinite form (since, by hypothesis, $F$
certainly is not a negative form), so it must be equivalent to the form
$xx-2yz$. We can therefore give a transformation of the latter into the
former, which provides a proper representation of the form $F'$
by $xx-2yz$. --- We will then have
\[
A=\alpha\alpha-2\alpha'\alpha'',\quad -B=\alpha\beta-\alpha'\beta''-\alpha''\beta',\quad C=\beta\beta-2\beta'\beta''
\]
\SourceLeaf{341}
Further, if we denote the numbers $\alpha\beta'-\alpha'\beta$, $\alpha'\beta''-\alpha''\beta'$, $\alpha''\beta-\alpha\beta''$ by $a,b,c$,
respectively, they will have no common divisor, and $D=bb-2ac$.

II.\ With the help of the last observation in Article 235, it follows
readily that $F$ is transformed by the substitution $2\beta',\allowbreak \beta,\allowbreak \beta,\allowbreak \beta'';\ 2\alpha',\allowbreak \alpha,\allowbreak \alpha,\allowbreak \alpha''$ into
the product of the form $(2a,-b,c)$ with itself, and by the substitution
$\beta',\allowbreak \beta,\allowbreak \beta,\allowbreak 2\beta'';\ \alpha',\allowbreak \alpha,\allowbreak \alpha,\allowbreak 2\alpha''$ into the product of the form $(a,-b,2c)$ with itself. Now the
greatest common divisor of the numbers $2a,2b,2c$ is $2$. So, if
$c$ is odd, $2a,2b,c$ will have no common divisor; that is, $(2a,-b,c)$ will
be a properly primitive form. Likewise, if $a$ is odd, $(a,-b,2c)$
will be a properly primitive form. In the first case $F$ arises
by doubling the form $(2a,-b,c)$, in the second by doubling the form
$(a,-b,2c)$ (see conclusion 4, Article 235); and one of these cases certainly
always occurs. For if both $a,c$ were even, $b$ would necessarily
be odd. We readily confirm that $\beta''a+\beta b+\beta'c=0$, $\alpha''a+\alpha b+\alpha'c=0$, which would imply
that $\beta b$, $\alpha b$, and therefore also $\alpha$ and $\beta$, were
even. But then $A$ and $C$ would be even, contrary to the
hypothesis that $F$ is a form of the principal genus and therefore
of the properly primitive order. --- Further, both $a$ and $c$
may be odd. In that case we immediately have two forms whose doubling
gives $F$.

\textit{Example.} Let the given form be $F=(5.\,2,31)$, of determinant $-151$.
A value of the expression $\surd(5,2,31)$ is found here to be $(55,22)$, giving
the ternary form $\varphi=\left(\begin{smallmatrix}5,&31,&4\\11,&0,&-2\end{smallmatrix}\right)$. The rules of Article 272 give an equivalent
form $\left(\begin{smallmatrix}1,&1.&-1\\0,&0,&0\end{smallmatrix}\right)$, which is transformed into $\varphi$ by the substitution
$\left\{\begin{smallmatrix}2,&2,&1\\1,&-6,&-2\\0,&3,&1\end{smallmatrix}\right\}$. Using the transformations given in Article 277, we then find
that $\left(\begin{smallmatrix}1,&0,&0\\-1,&0,&0\end{smallmatrix}\right)$ is transformed into $\varphi$ by the substitution $\left\{\begin{smallmatrix}3.&-7,&-2\\2,&-1,&0\\1,&-9,&-3\end{smallmatrix}\right\}$.
So $a=11$, $b=-17$, $c=20$. Since $a$ is odd, $F$ arises
by doubling the form $(11,17,40)$, and is transformed into the product of
this form with itself by the substitution $-1,-7,-7,-18;\ 2,3,3,2$.

